\documentclass{article}
\usepackage{amsmath,amssymb,xurl}
\usepackage[hidelinks]{hyperref}
% Shared commands for notes in docs/paper-gaps/.

\DeclareUnicodeCharacter{2080}{\ifmmode{}_{0}\else\textsubscript{0}\fi}
\DeclareUnicodeCharacter{2081}{\ifmmode{}_{1}\else\textsubscript{1}\fi}
\DeclareUnicodeCharacter{2082}{\ifmmode{}_{2}\else\textsubscript{2}\fi}
\DeclareUnicodeCharacter{2099}{\ifmmode{}_{n}\else\textsubscript{n}\fi}

\newcommand{\Lean}{\textsc{Lean}}
\ExplSyntaxOn
\NewDocumentCommand{\leanid}{m}{
  \ifmmode
    \textsf{\detokenize{#1}}
  \else
    \group_begin:
    \urlstyle{sf}
    \tl_set:Nn \l_tmpa_tl {#1}
    \tl_replace_all:Nnn \l_tmpa_tl {\_} {_}
    \exp_args:NV \path \l_tmpa_tl
    \group_end:
  \fi
}
\ExplSyntaxOff
\providecommand{\tr}{\operatorname{tr}}
\newcommand{\tnleanIssueBase}{https://github.com/LionSR/TNLean/issues/}
\newcommand{\tnleanPrBase}{https://github.com/LionSR/TNLean/pull/}
\newcommand{\tnleanDocBase}{https://sirui-lu.com/TNLean/docs/}
\newcommand{\ghissue}[1]{\href{\tnleanIssueBase#1}{GitHub issue \##1}}
\newcommand{\ghpr}[1]{\href{\tnleanPrBase#1}{PR \##1}}
\newcommand{\doclink}[1]{\href{\tnleanDocBase#1.html}{\path{#1}}}

% Dirac notation and the identity operator, as in blueprint/src/macros/common.tex.
% \providecommand keeps note-local definitions authoritative.
\usepackage{dsfont}
\providecommand{\ket}[1]{|#1\rangle}
\providecommand{\bra}[1]{\langle #1|}
\providecommand{\braket}[2]{\langle #1 | #2 \rangle}
\providecommand{\ketbra}[2]{|#1\rangle\!\langle #2|}
\providecommand{\Id}{\mathds{1}}
\providecommand{\one}{\mathds{1}}
\providecommand{\C}{\mathbb{C}}
\providecommand{\MN}[1]{M_{#1}(\C)}
\providecommand{\MD}{M_D(\C)}

% Verdict marker: \gapnote{<kind>}{<status>}. Kind is the policy
% classification (clarification, local-correction, scope-restriction,
% unfaithful, false-source, open-gap); status is open, wip (elimination
% underway), resolved, or historical. Machine-read by the site index;
% prints a verdict line.
\newcommand{\gapnote}[2]{\par\noindent\textbf{Verdict:} #1 (#2).\par}

\title{Finite Native Normal Forms for Proper Square Supports}
\author{}
\date{}
\begin{document}
\maketitle
\gapnote{scope-restriction}{open}

Schuch, Cirac, and P\'erez-Garc\'ia, \cite[Lemma~6.14]{Schuch2010PEPS},
source \path{Papers/1001.3807/paper_v3.tex}, lines 2199--2214,
deform a fixed-endpoint flux string by virtual invariance.
Here the allowed primal supports of the generated square moves are explicit.
This is a classification of that generated relation, not a statement that
all paths with common endpoints or displacement are equivalent.

Let $X=\mathbb Z_w\times\mathbb Z_h$ with $w,h>0$, and let $R\subsetneq X$
be the allowed primal supports. Write $p\sim_Rq$ for square moves supported
in $R$, backtracks, concatenation and reversal. The reverse-ordered native
word $W(p)$ lies in the free group on the positive east and north edges,
with each actual base vertex retained. The face with lower-left base $v$
has primal support $v+(1,1)$ and word $r_v=BAC^{-1}D^{-1}$, where $A,B,C,D$
are its bottom, right, left and top positive generators.
Its two substitutions are $C\mapsto D^{-1}BA$ and $A\mapsto B^{-1}DC$,
fixing every other generator.

There exists a finite list of those tagged substitutions, of length at most
$|X|$, using only faces supported in $R$, whose composition $\Sigma$ satisfies
\begin{align}
    \Sigma(r_v)=1\quad\text{if }v+(1,1)\in R.
    \label{eq:scp10_proper_support_face_elimination}
\end{align}
To construct it, choose $c\notin R$, put $o=c-(1,1)$, and rank each face base by
$\kappa(v)=\overline{v_1-o_1}+\overline{v_2-o_2}$, using standard coordinate
representatives. Repeatedly select a least-rank remaining allowed face.
A positive first coordinate selects west substitution; otherwise the second
coordinate is positive and selects south substitution. Its predecessor has
smaller rank, so the selected generator appears in no other remaining allowed
face. The required transverse period exceeds one because that shifted
coordinate is positive. Finite induction therefore kills each selected
relation without changing the others. This includes periods one and two.
\footnote{\leanid{TNLean.PEPS.exists_torusDualFaceEliminationSequence}.}

The same single list, constructed independently of the compared paths,
gives for all finite parallel native paths
\begin{align}
    p\sim_Rq
    \quad\Longleftrightarrow\quad
    \operatorname{red}(\Sigma(W(p)))
    =\operatorname{red}(\Sigma(W(q))).
    \label{eq:scp10_proper_support_normal_form}
\end{align}
Here $\operatorname{red}$ is the unique reduced signed word in the free group.
For necessity, left multiplication by each $\Sigma(W(e))$ gives a reversible
transport flat on every allowed face, so generated homotopies preserve
$\Sigma(W(p))$. For sufficiency, let $\pi_R$ be the explicit finite native
presentation quotient. Every selected face is an original defining relator.
Preservation of flat representations under finite substitutions gives
\begin{align}
    \pi_R\circ\Sigma=\pi_R.
    \label{eq:scp10_proper_support_quotient}
\end{align}
Reduced-word equality then implies equality of the original words in the
finite presentation, whose completeness supplies $p\sim_Rq$.
\footnote{\leanid{TNLean.PEPS.exists_torusDualProperSupportNormalForm}.}

This proves existence of a finite geometric normal form for every proper
support in the positive-period native model. It does not extract an executable
general Boolean decision procedure or give a complexity bound. Full support
retains the separate lifted-displacement classification; zero periods retain
the separate semantic theorem. An arbitrary missing face can still carry a
nontrivial nonabelian loop word, so no endpoint-only or displacement-only
classification follows. Physical density conclusions and annular Wilson-loop
measurements remain separate from this homotopy classification.
The broader source interpretation stays open under \ghissue{8703};
the Hamiltonian work under \ghissue{8947} is unchanged.

\bibliographystyle{alpha}
\bibliography{references}
\end{document}
